\documentclass[11pt,a4paper]{article}
\usepackage[margin=1in]{geometry}
\usepackage{amsmath,amssymb}
\usepackage{xcolor}
\usepackage[colorlinks=true,allcolors=blue]{hyperref}
\usepackage{enumitem}
\usepackage{titlesec}
\usepackage{parskip}
\usepackage{xr}
\externaldocument{build/main}

% Generated by make_numbers.py -- do not edit by hand.
\newcommand{\AblBaseline}{-0.0277}
\newcommand{\AblBestCell}{uniform / derived / all / tuned}
\newcommand{\AblBestRtwo}{0.7807}
\newcommand{\AblFull}{0.7622}
\newcommand{\AblNcells}{32}
\newcommand{\AblNorders}{24}
\newcommand{\CfgManualFull}{0.6925}
\newcommand{\CfgManualSub}{0.6901}
\newcommand{\CfgManualTopFull}{0.7438}
\newcommand{\CfgManualTopSub}{0.7155}
\newcommand{\CfgUniformFull}{0.7453}
\newcommand{\CfgUniformSub}{0.7423}
\newcommand{\CfgUniformTopFull}{0.7664}
\newcommand{\CfgUnscaledFull}{-0.0277}
\newcommand{\CfgUnscaledSub}{-0.0538}
\newcommand{\CvCorrhi}{0.719}
\newcommand{\CvCorrlo}{0.667}
\newcommand{\CvMean}{0.6931}
\newcommand{\CvSd}{0.0368}
\newcommand{\CvVOnehi}{0.765}
\newcommand{\CvVOnelo}{0.621}
\newcommand{\DSAmesBase}{0.6819}
\newcommand{\DSAmesFull}{0.8358}
\newcommand{\DSAmesFullUni}{0.8381}
\newcommand{\DSAmesN}{1022}
\newcommand{\DSAmesP}{36}
\newcommand{\DSAmesScal}{-0.1558}
\newcommand{\DSAmesSelect}{+0.1606}
\newcommand{\DSAmesTune}{+0.1491}
\newcommand{\DSAmesUniScal}{+0.1211}
\newcommand{\DSAmesUniSelect}{+0.0068}
\newcommand{\DSAmesUniTune}{+0.0283}
\newcommand{\DSCaliforniaBase}{0.5013}
\newcommand{\DSCaliforniaFull}{0.7563}
\newcommand{\DSCaliforniaFullUni}{0.7768}
\newcommand{\DSCaliforniaN}{10000}
\newcommand{\DSCaliforniaP}{18}
\newcommand{\DSCaliforniaScal}{+0.1616}
\newcommand{\DSCaliforniaSelect}{+0.0186}
\newcommand{\DSCaliforniaTune}{+0.0749}
\newcommand{\DSCaliforniaUniScal}{+0.1869}
\newcommand{\DSCaliforniaUniSelect}{+0.0157}
\newcommand{\DSCaliforniaUniTune}{+0.0730}
\newcommand{\DSKcHouseBase}{0.3888}
\newcommand{\DSKcHouseFull}{0.8092}
\newcommand{\DSKcHouseFullUni}{0.7806}
\newcommand{\DSKcHouseN}{10000}
\newcommand{\DSKcHouseP}{20}
\newcommand{\DSKcHouseScal}{+0.0805}
\newcommand{\DSKcHouseSelect}{+0.1197}
\newcommand{\DSKcHouseTune}{+0.2202}
\newcommand{\DSKcHouseUniScal}{+0.2710}
\newcommand{\DSKcHouseUniSelect}{+0.0053}
\newcommand{\DSKcHouseUniTune}{+0.1155}
\newcommand{\FlipA}{SVR uniform / raw8 / default / n=3000}
\newcommand{\FlipB}{SVR manual / top12 / tuned / full}
\newcommand{\FlipDiff}{-0.0271}
\newcommand{\FlipP}{0.166}
\newcommand{\FlipPuncorr}{0.007}
\newcommand{\GbmFull}{0.8384}
\newcommand{\GbmSmall}{0.8055}
\newcommand{\GridFits}{192}
\newcommand{\GridRtwo}{0.7806}
\newcommand{\GridSecs}{422}
\newcommand{\GridSecsMin}{7.0}
\newcommand{\KbelowPlateau}{0.5745}
\newcommand{\KbestAt}{10}
\newcommand{\KbestK}{10}
\newcommand{\KernLinear}{0.6469}
\newcommand{\KernPolyThree}{0.6650}
\newcommand{\KernPolyTwo}{0.6613}
\newcommand{\KernRbf}{0.7702}
\newcommand{\KernSigmoid}{0.3207}
\newcommand{\KmaxAt}{18}
\newcommand{\KminAt}{4}
\newcommand{\KnnFull}{0.7599}
\newcommand{\KnnSmall}{0.7179}
\newcommand{\KplateauFrom}{10}
\newcommand{\KplateauRange}{0.0045}
\newcommand{\KrangeHi}{0.7683}
\newcommand{\KrangeLo}{0.5745}
\newcommand{\MaxDeriv}{+0.5428}
\newcommand{\MaxScal}{+0.7202}
\newcommand{\MaxSelect}{+0.5937}
\newcommand{\MaxTune}{+0.5312}
\newcommand{\MinDeriv}{-0.0881}
\newcommand{\MinScal}{-0.0239}
\newcommand{\MinSelect}{-0.0029}
\newcommand{\MinTune}{+0.0152}
\newcommand{\NDatasets}{3}
\newcommand{\NFlipped}{2}
\newcommand{\NModels}{10}
\newcommand{\NOrderHoldsAuto}{1}
\newcommand{\NOrderHoldsUni}{3}
\newcommand{\NPairs}{21}
\newcommand{\NestedRandom}{0.7532}
\newcommand{\NestedRandomMax}{0.767}
\newcommand{\NestedRandomMin}{0.744}
\newcommand{\NestedSpatial}{-0.5344}
\newcommand{\NestedSpatialMax}{0.660}
\newcommand{\NestedSpatialMin}{-4.091}
\newcommand{\NystroemM}{2000}
\newcommand{\NystroemRtwo}{0.7673}
\newcommand{\NystroemSecs}{18}
\newcommand{\PairManualDiff}{+0.0242}
\newcommand{\PairManualP}{4.6\times 10^{-6}}
\newcommand{\PairManualPuncorr}{3.4\times 10^{-9}}
\newcommand{\PairRfDiff}{-0.0446}
\newcommand{\PairRfP}{9.3\times 10^{-5}}
\newcommand{\PairRfPuncorr}{9.4\times 10^{-8}}
\newcommand{\PairXgbDiff}{-0.0728}
\newcommand{\PairXgbP}{1.9\times 10^{-7}}
\newcommand{\PairXgbPuncorr}{1.2\times 10^{-10}}
\newcommand{\PreethiReportedMSE}{1.12}
\newcommand{\PreethiReportedRtwo}{0.15}
\newcommand{\RandSixtyRtwo}{0.7834}
\newcommand{\RandSixtySecs}{203}
\newcommand{\RandTwentyRtwo}{0.7702}
\newcommand{\RandTwentySecs}{100}
\newcommand{\RandTwoHundredRtwo}{0.7796}
\newcommand{\RandTwoHundredSecs}{1377}
\newcommand{\RandomCVmean}{0.7496}
\newcommand{\RangeDeriv}{0.6309}
\newcommand{\RangeScal}{0.7441}
\newcommand{\RangeSelect}{0.5966}
\newcommand{\RangeTune}{0.5160}
\newcommand{\RepScaledDefault}{0.737}
\newcommand{\RepScaledTuned}{0.766}
\newcommand{\RepScaledTunedMSE}{0.309}
\newcommand{\RepSeeds}{5}
\newcommand{\RepUnscaledDefault}{-0.020}
\newcommand{\RepUnscaledDefaultMSE}{1.343}
\newcommand{\RepUnscaledTuned}{0.514}
\newcommand{\RepUnscaledTunedMSE}{0.640}
\newcommand{\RfFull}{0.8171}
\newcommand{\RfSmall}{0.7610}
\newcommand{\RidgeFull}{0.6504}
\newcommand{\RidgeSmall}{0.6497}
\newcommand{\SearchBest}{random 60}
\newcommand{\SearchBestRtwo}{0.7834}
\newcommand{\SearchBestSecs}{203}
\newcommand{\SelectionOptimism}{0.0000}
\newcommand{\ShapDeriv}{0.1505}
\newcommand{\ShapScal}{0.2723}
\newcommand{\ShapSelect}{0.1789}
\newcommand{\ShapTune}{0.1882}
\newcommand{\SizeFull}{0.7813}
\newcommand{\SizeFullN}{14448}
\newcommand{\SizeFullRmse}{0.5382}
\newcommand{\SizeFullSecs}{8.3}
\newcommand{\SizeSmall}{0.6775}
\newcommand{\SizeSmallN}{1000}
\newcommand{\SizeThreeK}{0.7418}
\newcommand{\SizeThreeKrmse}{0.5848}
\newcommand{\SpatialCVmean}{0.2316}
\newcommand{\SpatialCVsd}{0.7525}
\newcommand{\SpatialOptimism}{0.5180}
\newcommand{\StatRho}{0.429}
\newcommand{\StatSplits}{10}
\newcommand{\StatSvrCIhi}{0.783}
\newcommand{\StatSvrCIlo}{0.769}
\newcommand{\StatSvrMean}{0.7757}
\newcommand{\StatSvrNBhi}{0.792}
\newcommand{\StatSvrNBlo}{0.759}
\newcommand{\StatSvrSd}{0.0099}
\newcommand{\StatSvrVOnehi}{0.795}
\newcommand{\StatSvrVOnelo}{0.756}
\newcommand{\SvrRankFull}{4}
\newcommand{\SvrRankSmall}{4}
\newcommand{\SvrRbfFull}{0.7702}
\newcommand{\SvrRbfSmall}{0.7501}
\newcommand{\ThreshRange}{0.0004}
\newcommand{\WeightMax}{0.7674}
\newcommand{\WeightMin}{0.7128}
\newcommand{\WeightMinOverlap}{10}
\newcommand{\WeightN}{66}
\newcommand{\WeightRange}{0.0546}
\newcommand{\XgbFull}{0.8483}
\newcommand{\XgbSmall}{0.8057}


\definecolor{revcol}{gray}{0.32}
\newcommand{\comment}[1]{\par\smallskip{\color{revcol}\itshape #1}\par\smallskip}
\newcommand{\reply}{\par\noindent\textbf{Response.}\ }
\newcommand{\pt}[1]{\par\medskip\noindent\textbf{#1}\par}

\titleformat{\section}{\large\bfseries}{\thesection.}{0.6em}{}
\setlength{\emergencystretch}{2em}

\begin{document}

\begin{center}
{\Large\bfseries Response to Reviewers}\\[4pt]
{\large A Controlled Reassessment of Support Vector Regression on the California\\
Housing Benchmark with Feature Scaling, Feature Selection and Hyperparameter Tuning}\\[6pt]
Emmanuel Adutwum \quad$\cdot$\quad Soka University of America \quad$\cdot$\quad \texttt{eadutwum@soka.edu}\\[4pt]
\emph{Discover Artificial Intelligence} \quad$\cdot$\quad Submission ID 6cf4e30e-1a5b-4ad2-99e5-2320f0737507
\end{center}

\vspace{4pt}
\hrule
\vspace{10pt}

\noindent Dear Dr Addula and Mr Mohite,

Thank you for the decision of major revision, and my thanks to the three reviewers for
reports that were detailed, technically specific and, in two important places, correct
about errors I had made. The revision that follows is substantial: the manuscript has been
rewritten around a new set of experiments, and two of the earlier version's central claims
are withdrawn rather than defended.

Three changes deserve to be stated at the outset, because they alter the paper's central
argument rather than merely elaborating it.

\pt{1. The reference value was misread, and the correction removes the paper's central
inconsistency.}
Reviewer~1 (point~2) and Reviewer~3 both observed that an unscaled baseline of
$R^2 = -0.054$ cannot explain a previously reported $R^2$ of $0.60$, and that the earlier
manuscript never accounted for the gap. They were right, and on re-reading Preethi et al.\
closely I found the cause: \textbf{$0.60$ is not their SVR result.} It is the score of their
linear, ridge and polynomial-ridge models. Their SVR-RBF model attains
$R^2 \approx \PreethiReportedRtwo$ with MSE $\approx \PreethiReportedMSE$, read from their
Figs.~7 and~8 (their paper reports no numerical table). The earlier version of my manuscript
attributed the linear family's score to their SVR. That was my error. With the correct value
in place, the configuration they describe reproduces at $R^2 = \RepUnscaledDefault$ with
library defaults and $\RepUnscaledTuned$ when $C$ and $\varepsilon$ are tuned --- values that
bracket $\PreethiReportedRtwo$ --- and no unexplained gap remains. This is
Section~\ref{sec:the-result} and Section~\ref{sec:res-reproduction} of the revised manuscript.

\pt{2. The proposed feature-specific scaling strategy does not work, and is now reported as
a negative result.}
Reproducing the prior study's uniform scaling directly, as the Editor asked in point~1,
required me to compare it against my own feature-specific assignment on an identical split.
Uniform standardisation gives $R^2 = \CfgUniformFull$; the hand-assigned feature-specific
strategy gives $\CfgManualFull$ on the same split with the same features. The strategy I
presented in the earlier version as the preprocessing advance is therefore
\emph{worse} than plain standardisation. It is now reported as such, and the workflow's
scaling stage is described as ``any reasonable scaling'' rather than as a contribution.

\pt{3. The attribution of 95.7\% of the gain to scaling was an artefact of ablation order.}
Evaluating the full factorial design rather than a single path shows that the apparent
contribution of scaling ranges over \RangeScal{} in $R^2$ depending only on where it sits in
the sequence. The single-path figure reported earlier sat near the top of that range. The
order-independent Shapley value is $\ShapScal$. Scaling remains the largest contributor,
but the earlier percentage is withdrawn.

I have also added a disclosure of my use of a large language model, as Reviewer~2 requested
and the Editor endorsed in point~12.

Every point raised by the Editor and all three reviewers is answered below, with the
location of each change. Reviewer comments are set in \textcolor{revcol}{\itshape grey
italic}. All numerical values quoted in this letter are generated from the same result
records that generate the manuscript's tables and figures, so the two documents cannot
disagree.

\vspace{6pt}
\hrule

\section{Response to the Editor's twelve points}

\pt{E1. Directly reproduce baseline performance.}
\comment{The uniform scaling SVR configuration used in previous work should be reproduced
directly and its performance compared to the authors' feature-wise scaling method.}
\reply Done, and it produced a negative result that I report. A new experiment
(Section~\ref{sec:res-reproduction}, Tables~2 and~3) re-runs the configuration described by Preethi et al.\ over
\RepSeeds{} seeds: eight raw features, RBF kernel, a grid over $C$ and $\varepsilon$ only,
$\gamma$ at the library default, 80/20 partition. It yields $R^2 = \RepUnscaledDefault$
untuned and $\RepUnscaledTuned$ tuned, bracketing their reported
$\PreethiReportedRtwo$. Table~\ref{tab:configurations} then compares uniform standardisation against my
feature-specific assignment on one identical split across all feature sets and training
sizes: uniform gives $\CfgUniformFull$, feature-specific gives $\CfgManualFull$. The
feature-specific strategy loses. This is stated in the abstract, in
Section~\ref{sec:res-reproduction}, and again in Section~\ref{sec:disc-not} among the withdrawn claims.

\pt{E2. Fair model comparison.}
\comment{SVR and baseline algorithms should be trained / evaluated with similar quantities
of data, rather than trained-set sizes that differ by an order of magnitude or more.}
\reply Corrected. Section~\ref{sec:res-models} (Table~\ref{tab:models}) retrains all \NModels{} models with identical
features, identical training rows and an identical search budget of 20 randomised iterations
with three-fold cross-validation, reported at both 3{,}000 rows and the full training set.
The tuned SVR attains $\SvrRbfSmall$ (rank \SvrRankSmall) at 3{,}000 rows and $\SvrRbfFull$
(rank \SvrRankFull) on the full set. Section~\ref{sec:res-kernels} (Table~\ref{tab:size}) additionally shows that the exact
full-data fit takes only $\SizeFullSecs$~s, so the subset restriction had no computational
justification and has been removed from the main results rather than defended. A Nystr\"om
approximation is reported for completeness but is not needed at this scale.

\pt{E3. Statistical significance testing.}
\comment{The confidence interval equation is incorrect. Proper paired statistical tests
should be performed to compare models.}
\reply Both parts corrected. The earlier version reported $\bar{r} \pm 1.96 s$, which is a
normal-approximation prediction interval for a single fold, not a confidence interval for
the mean. Section~\ref{sec:res-stats} (Table~\ref{tab:statistics}) now reports the interval three ways so the difference is
explicit: the correct $\bar{r} \pm t_{0.975,J-1}\, s/\sqrt{J}$ gives
$[\StatSvrCIlo, \StatSvrCIhi]$, the Nadeau--Bengio interval gives
$[\StatSvrNBlo, \StatSvrNBhi]$, and the earlier expression gives
$[\StatSvrVOnelo, \StatSvrVOnehi]$. The ten-fold interval of the earlier version moves from
$[\CvVOnelo, \CvVOnehi]$ to $[\CvCorrlo, \CvCorrhi]$. The argument that the earlier interval
did not overlap a competing value is withdrawn. For model comparison, all models are now
evaluated on \StatSplits{} \emph{identical} splits and compared with the Nadeau--Bengio
corrected resampled $t$-test, since resampled training sets overlap and the ordinary paired
$t$ is anti-conservative; the uncorrected $t$ and the Wilcoxon signed-rank test are reported
alongside it (Table~\ref{tab:paired}).

\pt{E4. Ablation study should be improved.}
\comment{The impact of each stage of the proposed pipeline should be tested more
rigorously, including the order in which they are applied or alternative methods to several
pipeline stages.}
\reply Rebuilt. The ablation is now a four-factor full factorial design --- scaling, derived
features, selection, tuning --- and \emph{every one} of its \AblNcells{} cells is evaluated
(Table~\ref{tab:ablation-cells}). Each component receives a Shapley value, which averages over all \AblNorders{}
orderings and is the unique attribution that is both order-independent and additive. The
result changed the paper's conclusion: scaling appears to contribute between $\MinScal$ and
$\MaxScal$ depending on placement, and its order-independent value is $\ShapScal$, not the
95.7\% previously claimed. Alternative methods at each stage are also swept: three scaling
strategies, four kernels each tuned over its own hyperparameters (Table~\ref{tab:kernels}), and grid versus
randomised search at three budgets (Table~\ref{tab:search}).

\pt{E5. Validate predictive inference.}
\comment{Consider spatial validation of the model, or at minimum holdout feature selection
within each cross-validation fold to avoid overly optimistic estimates of model
performance.}
\reply Both were done, and this produced the most consequential result in the paper.
Section~\ref{sec:res-validation} (Table~\ref{tab:validation}) crosses two partitioning schemes with two selection schemes. Random
ten-fold cross-validation gives $\RandomCVmean$; cross-validation over ten spatially
contiguous blocks, obtained by $k$-means on standardised centroids, gives
$\SpatialCVmean \pm \SpatialCVsd$. The optimism of random partitioning, $\SpatialOptimism$,
is larger than every pipeline effect measured anywhere in this study. Feature selection
repeated inside each fold changes the estimate by $\SelectionOptimism$; the effect is small
because the selected sets are stable across folds, but it is reported because the earlier
version could not have known that. Fully nested estimates, with both selection and
hyperparameter search inside each outer fold, are $\NestedRandom$ (random) and
$\NestedSpatial$ (spatial). Section~\ref{sec:disc-threats} now states plainly that results of this kind measure
interpolation within a sampled region, not prediction for an unsampled one.

\pt{E6. Limit claims to dataset and configuration studied.}
\comment{Given that only one benchmark dataset was used, claims in the discussion should be
limited to those pertaining to the California Housing dataset and specific configuration
tested.}
\reply Done throughout. A new subsection, \emph{Scope of the claims} (Section~\ref{sec:scope}), states
that all quantitative claims are restricted to the datasets, feature sets, search spaces and
implementations described in Section~\ref{sec:setup}, and that the necessity of scaling for distance-based
kernels is long established and not claimed here as a discovery. The title now reads ``A
Controlled Reassessment \ldots on the California Housing Benchmark''. Section~\ref{sec:disc-not} is devoted
to what the evidence does \emph{not} support. The abstract closes on the same restriction.
Every general statement about SVR has been removed or bounded.

\pt{E7. Define proposed workflow.}
\comment{The overall processing pipeline should be explicitly defined. It may also be
helpful to include a diagram that illustrates the different stages of the pipeline and what
data is input and output to each stage.}
\reply A new Section~\ref{sec:workflow}, \emph{Proposed workflow}, defines the pipeline as nine stages
(S1--S9) with stated inputs and outputs, \emph{before} any data are analysed. Figure~\ref{fig:pipeline} is a
new diagram showing each stage and the data flowing between them, with the stages fitted on
the training fold enclosed in a dashed boundary; Table~\ref{tab:workflow} states the same information
tabularly. The scaler assignment, previously a hand-made table specific to one dataset, is
now stated as an algorithm (Algorithm~\ref{alg:scaler}) that assigns a scaler to each feature from
training-fold statistics alone, so it can be audited and applied elsewhere.

\pt{E8. Definition of methods is lacking.}
\comment{Much more detail is needed in the description of how models are trained, how
hyperparameters were optimized, what cross-validation strategy was used, how features were
engineered and models compared, and how the ablation study was performed.}
\reply Section~\ref{sec:workflow} now specifies each of these separately: Section~\ref{sec:svr} the SVR formulation
with every symbol defined; Sections~\ref{sec:derived}--\ref{sec:scaler} feature derivation, ensemble importance and
selection; Section~\ref{sec:tuning} the scaler rule as pseudocode; Section~\ref{sec:eval-protocol} the search spaces and
budgets; Section~\ref{sec:ablation-protocol} the evaluation protocol, which distinguishes explicitly between the
inner cross-validation used for hyperparameter selection and the outer resampling used for
performance estimation --- the ambiguity Reviewer~2 identified; and Section~\ref{sec:ablation-protocol} the ablation
protocol. Section~\ref{sec:fixed}, \emph{What is held fixed}, states for each comparison exactly which
factor varies and which are held constant.

\pt{E9. Claim of novelty should be improved.}
\comment{Currently the paper suggests that the authors have broadly ``rediscovered'' SVR.
The authors should avoid making this claim unless a new novel contribution is demonstrated.}
\reply Withdrawn entirely. Section~\ref{sec:contributions} now opens: ``The contribution of this paper is
methodological and diagnostic rather than algorithmic. No new learning algorithm is
proposed.'' Section~\ref{sec:rw-scaling} places the work explicitly within the existing scaling literature
(Sola and Sevilla; Hsu et al.; Singh and Singh; Fern\'andez-Delgado et al.) and states that
the paper does not claim that result. What is claimed is narrower and, I hope, defensible: a
corrected reading of one published result, a direct reproduction of it, an order-independent
ablation procedure, a corrected evaluation protocol, and two negative results.

\pt{E10. Manuscript language and formatting requires revision.}
\comment{Figure / table captions, references, math notation, numbering of results all
require revision. In addition, the overall flow of the manuscript can be improved by
revising section headings.}
\reply The manuscript was rewritten rather than patched. Specifically: figure titles
embedded in images have been removed so the caption is the sole description; all tables are
referenced by arabic numeral matching their captions, replacing the Roman numerals of the
earlier version; every reference has been rechecked and DOIs added; all mathematical symbols
are defined where introduced (Section~\ref{sec:svr}) and collected in a notation table (Table~\ref{tab:notation});
section headings are now declarative statements of content, and Section~\ref{sec:workflow} (workflow) now
precedes Section~\ref{sec:setup} (experimental setup) as Reviewer~2 requested. On numerical consistency I
took a structural approach rather than a proofreading one: \textbf{no number is typed by hand
anywhere in this manuscript}. Every experiment writes a JSON record; a generator script turns
those records into every table, every figure and every in-text value. The class of error
Reviewer~1 identified in point~9 is therefore now impossible rather than merely corrected.

\pt{E11. Use additional datasets.}
\comment{Test your proposed workflow on additional datasets. Reviewer \#2 particularly felt
that this would improve the authors' argument that their method does not only work well for
California Housing.}
\reply Section~\ref{sec:res-datasets} (Table~\ref{tab:datasets}) applies the workflow unchanged --- same scaler rule, same
selection rule, same search space, no dataset-specific adjustment --- to King County house
sales (21{,}613 transactions, 18 numerical predictors) and Ames Housing (1{,}460 sales, 36
usable numerical predictors). The result is not the clean confirmation I expected, and I
report it as it came out.

Under plain uniform standardisation the ordering scaling $>$ tuning $>$ selection is
reproduced on all \NOrderHoldsUni{} datasets: scaling contributes $\DSCaliforniaUniScal$,
$\DSKcHouseUniScal$ and $\DSAmesUniScal$, tuning $\DSCaliforniaUniTune$,
$\DSKcHouseUniTune$ and $\DSAmesUniTune$, and selection $\DSCaliforniaUniSelect$,
$\DSKcHouseUniSelect$ and $\DSAmesUniSelect$.

Under the distribution-aware scaler rule that this revision introduced, it is reproduced on
only \NOrderHoldsAuto{} of \NDatasets{}. On King County the rule makes scaling the smallest
contributor ($\DSKcHouseScal$, against $\DSKcHouseTune$ for tuning); on Ames its Shapley
value is negative ($\DSAmesScal$), and applying it first costs $0.6455$ in $R^2$. Ames is the
smallest and widest of the three ($\DSAmesN$ training rows, $\DSAmesP$ candidate features),
which is where per-feature scaler assignment estimated from training statistics is least
reliable.

I have therefore withdrawn the rule as a contribution rather than presenting the additional
datasets as a confirmation. Section~\ref{sec:res-datasets} reports both strategies side by
side, Figure~\ref{fig:datasets} plots both, and Section~\ref{sec:disc-not} lists the rule
among the claims the evidence does not support. This is the third of the paper's negative
results, and it is the one that most directly answers Reviewer~2's concern: testing the
workflow on other datasets did not corroborate it, it falsified part of it. I also state the
limit of the evidence in Section~\ref{sec:disc-threats}: all three are residential-property
datasets, so this tests transfer across datasets within a domain, not across domains.

\pt{E12. AI-use statement.}
\comment{Reviewer \#2 found large portions of the text to be generated by AI. Authors should
consider adding a statement disclosing use of any LLM.}
\reply Added as a declaration subsection, \emph{Use of generative artificial intelligence}.
It discloses that I used a large language model (Anthropic Claude) as an assistive tool for
drafting and copy-editing prose, for implementing and refactoring the experiment scripts, and
for generating tables and figures from the result records; and it states that the research
questions, experimental design, interpretation and all scientific claims are my own, that I
have reviewed and verified every number, statement and reference, and that no text appears
that I have not read and checked against the underlying results. Separately, the specific
symptoms Reviewer~2 identified --- an excessive number of thin subsections, and mismatches
between figures and captions --- have been fixed by consolidating the subsections and
regenerating every figure with its caption from the same source.

\section{Response to the internal editorial feedback}

\pt{I1. Declaration section with individual subheadings.}
\reply The Declarations section now carries \emph{Ethical approval}, \emph{Consent to
participate} and \emph{Consent to publish} as separate subheadings, each answered
individually. All three are ``Not applicable'', with the reason stated: the study analyses
publicly available aggregated datasets and involves no human participants, animal subjects,
or identifiable personal data.

\pt{I2. Data availability statement in the manuscript above the References.}
\reply Added above the References, with direct links: the California Housing dataset via
\texttt{sklearn.datasets.fetch\_california\_housing} and the StatLib archive; King County via
OpenML \url{https://www.openml.org/d/42731}; Ames via OpenML
\url{https://www.openml.org/d/42165}. The analysis code, the JSON records of every experiment
reported, and the scripts that generate every table and figure are at
\url{https://github.com/emmanueladutwum123/svr-california-housing}.

\pt{I3. Dual publication declaration.}
\reply Added. An earlier version was deposited as a preprint on arXiv (arXiv:2605.08660);
no part of the work is published in or under consideration at any other peer-reviewed
journal. The declaration states that the present manuscript is a substantially revised
version of that preprint, that the reading of the reference result is corrected, and that the
experimental content of Sections~\ref{sec:res-reproduction}--\ref{sec:res-kernels} is new.

\section{Response to Reviewer 1}

I am grateful for this report. Points~2 and~9 identified a genuine error in the earlier
manuscript, and correcting it changed the paper.

\pt{R1.1 Reframe as a controlled reassessment, not a new general finding.}
\comment{The improvement of SVR performance through feature scaling is a long-established
result \ldots\ The manuscript should therefore be reframed not as presenting a new general
finding about SVR, but as a controlled reassessment of the specific result reported by
Preethi et al.\ (2025).}
\reply Accepted in full; this is now the frame of the paper, beginning with its title.
Section~\ref{sec:scope} states that scaling for distance-based kernels is long established and cites
Sola and Sevilla (1997), Hsu et al.\ (2003) and Singh and Singh (2020) for it; Section~\ref{sec:rw-scaling}
positions the work within that literature; Section~\ref{sec:disc-not} repeats that the paper establishes
nothing about SVR in general.

\pt{R1.2 The unscaled baseline cannot explain the previously reported 0.60.}
\reply You were right, and finding out why changed the paper. The $0.60$ figure is the score
of their \emph{linear family}; their SVR is $R^2 \approx \PreethiReportedRtwo$. My earlier
manuscript misattributed it. Section~\ref{sec:res-reproduction} now reproduces their configuration directly, under
the same 80/20 split over \RepSeeds{} seeds, obtaining $\RepUnscaledDefault$ untuned and
$\RepUnscaledTuned$ tuned, which bracket their value; and Table~\ref{tab:configurations} gives the direct comparison
of their uniform scaling ($\CfgUniformFull$) against my feature-specific strategy
($\CfgManualFull$) on an identical split. Thank you for pressing on this.

\pt{R1.3 SVR trained on 3{,}000 rows against competitors trained on 14{,}448.}
\reply Corrected; see E2 above. All \NModels{} models now receive identical rows, features
and search budget (Table~\ref{tab:models}), reported at both sizes. Table~\ref{tab:size} shows the full-data fit costs
$\SizeFullSecs$~s, so the subset had no justification. Nystr\"om approximation is reported
but proves unnecessary at this scale.

\pt{R1.4 Paired tests on identical splits rather than a bare confidence interval.}
\reply Adopted, using both references you gave. Section~\ref{sec:res-stats} evaluates every model on
\StatSplits{} identical splits and reports paired differences with Nadeau--Bengio (2003)
corrected intervals and $p$-values, with the uncorrected $t$ and Wilcoxon alongside for
reference (Dietterich, 1998). Against XGBoost the difference is $\PairXgbDiff$ with corrected
$p = \PairXgbP$, where the uncorrected test would have given $p = \PairXgbPuncorr$ --- I
report both precisely because the discrepancy makes your point.

\pt{R1.5 The ablation follows only one sequence.}
\reply Corrected; see E4. All \AblNcells{} cells of the four-factor design are evaluated and
attribution is by Shapley value over all \AblNorders{} orderings. The order-dependence you
anticipated is large: scaling's apparent contribution spans \RangeScal{} in $R^2$. The
manuscript now states explicitly that single-path increments should be read as upper bounds
for whichever component was applied first.

\pt{R1.6 Spatially blocked validation.}
\reply Added, following Roberts et al.\ (2017); see E5. Spatial blocking costs
$\SpatialOptimism$ in $R^2$ relative to random partitioning --- more than any pipeline effect
in the paper. Figure~\ref{fig:validation} shows the blocks and the resulting estimates.

\pt{R1.7 Justify the engineered features, the 40/30/30 weights and the top-12 selection.}
\comment{In particular, Coastal Proximity is based only on latitude difference, while
Location Score is defined through the product of latitude and longitude \ldots\ preferably
through sensitivity analysis.}
\reply Two changes. First, you were right that neither name was defensible, and
Section~\ref{sec:derived} now says so in the text rather than leaving the labels to imply
more than they measure. \emph{Coastal proximity}, $|\text{Latitude} - 34.05|$, is stated to
be a one-dimensional north--south displacement from a reference latitude near the Los Angeles
coastline --- explicitly \emph{not} a distance to the coast, which would require the
coastline geometry and both coordinates. \emph{Location score},
$(\text{Latitude} \times \text{Longitude})/1000$, is stated to be an interaction term between
the two coordinates with no geographic units, to be read as an unnamed coordinate interaction
rather than a location measure. Both are retained as \emph{candidate} features whose value is
decided empirically at Stage~S4 rather than by their names, and the ablation reports how
little the derived block contributes once selection and tuning are present. Second, Section~\ref{sec:res-sensitivity} (Table~\ref{tab:sensitivity}) sweeps
every hand-chosen constant: across \WeightN{} weight combinations on the simplex test $R^2$
varies only between \WeightMin{} and \WeightMax{} (range \WeightRange{}), and the smallest
overlap between any swept selection and the default top-twelve is \WeightMinOverlap{} of
twelve; and varying the scaler thresholds over nine combinations moves it by \ThreshRange{}.
Those two constants are therefore not load-bearing.

The number of selected features is a different matter, and the sweep changed what I can
claim. Varying $k$ from \KminAt{} to \KmaxAt{} moves $R^2$ between \KrangeLo{} and
\KrangeHi{} --- a spread of almost $0.2$, larger than any component effect in the ablation.
The dependence is not smooth: $R^2$ rises steeply to $k = \KplateauFrom$ and is then flat,
varying by only \KplateauRange{} across all $k$ from \KplateauFrom{} to \KmaxAt{}. So the
defensible statement is not that $k$ does not matter, but that it matters until about ten
features are retained and not afterwards, and that the default of twelve sits on the
plateau. I have written it that way in Section~\ref{sec:res-sensitivity} rather than
reporting the range alone, because the range alone would have implied a robustness the data
do not show.

\pt{R1.8 Feature selection performed outside the cross-validation folds.}
\reply Corrected rather than merely acknowledged. Selection is now repeated inside every
fold, and Table~\ref{tab:validation} reports selection-inside against selection-outside under both partitioning
schemes; the difference is $\SelectionOptimism$. Fully nested estimates following Varma and
Simon (2006) are $\NestedRandom$ (random) and $\NestedSpatial$ (spatial). The effect is small
here because the selected sets are stable across folds --- feature-selection stability is now
recorded in the result files --- but the earlier version had no way of knowing that, which
was your point.

\pt{R1.9 Findings rest on a single benchmark; restrict conclusions and moderate broad
statements throughout the title, abstract, discussion and conclusion.}
\comment{The findings are based on a single, extensively studied benchmark dataset. Unless
additional datasets are included, all conclusions should be restricted to the California
Housing dataset and the specific experimental configuration used. Broad statements regarding
the general performance of SVR should be moderated throughout the title, abstract, discussion,
and conclusion.}
\reply Addressed on both of the paths you offer, since I took the first one as well as the
second. Additional datasets \emph{are} now included: Section~\ref{sec:res-datasets} applies
the workflow unchanged to King County house sales and Ames Housing, and reports that the
ordering of the components is reproduced on both under uniform standardisation, while the
distribution-aware scaler rule introduced in this revision fails to transfer and is withdrawn
as a contribution. Independently
of that, the claims have been moderated in each of the four places you name. The
\textbf{title} is now ``A Controlled Reassessment \ldots on the California Housing
Benchmark''. The \textbf{abstract} closes on the sentence ``Conclusions are restricted to the
datasets and configurations examined here.'' The \textbf{discussion} opens each of its
supported claims with ``Within the California Housing dataset and the configurations tested
here'', and Section~\ref{sec:disc-not} states that the paper establishes nothing about SVR in
general --- only that one published comparison was configured in a way that disadvantaged one
of the models it compared. The \textbf{conclusion} restates the restriction explicitly. A new
Section~\ref{sec:scope} states the scope of every quantitative claim before any result is
presented. I would also note the limit of the new evidence rather than overstate it: all
three datasets are residential-property datasets, so Section~\ref{sec:res-datasets} tests
transfer within a domain, not across domains, and Section~\ref{sec:disc-threats} says so.

\pt{R1.10 Numerical and reporting inconsistencies throughout.}
\reply Every number in the manuscript has been recalculated from scratch, and the class of
error has been eliminated structurally: each experiment writes a JSON record and a generator
produces every table, figure and in-text value from those records, so no value is typed by
hand and the text, tables and figures cannot diverge. Rounding is standardised to four
decimal places for $R^2$ and to the reported precision elsewhere. The monetary conversions
you flagged have been removed rather than corrected, since with the target capped at
\$500{,}001 they were more misleading than informative. Mean and median are now labelled
distinctly wherever both appear. Figure and table numbering is generated automatically.

\section{Response to Reviewer 2}

This report identified the two structural weaknesses I found hardest to answer --- that the
pipeline was never actually formalised, and that the manuscript showed signs of
over-generated text. Both are addressed directly.

\pt{R2.0 (general remark) Figures 3--6 are unnecessary for a well-known dataset, or at
least not at their current size.}
\reply Agreed and acted on: the exploratory figures of the earlier version --- the
distribution and correlation displays that occupied Figures 3--6 --- have been removed from
the manuscript entirely rather than merely reduced. The dataset is, as you say, thoroughly
documented elsewhere, and Section~\ref{sec:datasets} now conveys the two facts that actually
matter for this paper's argument in a sentence of text: the feature ranges that motivate the
scaling stage, and the census recording cap at \$500{,}001 that bounds achievable fit. Every
figure that remains reports a result rather than describing the data.

\subsection*{Major issues}

\pt{R2.1 The contribution is weak; the pipeline is explained vaguely and should be
formalised with a diagram.}
\reply Accepted. Section~\ref{sec:contributions} now states plainly that no new learning algorithm is proposed
and that the contribution is methodological and diagnostic. The pipeline is formalised in a
new Section~\ref{sec:workflow} as stages S1--S9 with defined inputs and outputs, with Figure~\ref{fig:pipeline} as the diagram
you suggested and Table~\ref{tab:workflow} as its tabular statement, and the scaler assignment promoted from a
hand-made table to an auditable algorithm (Algorithm~\ref{alg:scaler}). I would add that the revision now
also offers a methodological contribution that the earlier version did not: the
order-independent Shapley attribution for preprocessing ablations (Section~\ref{sec:ablation-protocol}), which is
what revealed that the earlier 95.7\% figure was an artefact.

\pt{R2.2 Why compare only against Preethi et al.? Why only four references in Section~\ref{sec:rw-benchmark}?}
\reply Section~\ref{sec:related} has been expanded and restructured into three subsections --- scaling in
kernel methods, the California Housing benchmark, and model comparison and attribution ---
with the benchmarking literature discussed in terms of settings and reported performance
rather than merely listed. Preethi et al.\ remain the specific object of the reassessment,
because a controlled reproduction requires a specific configuration to reproduce, but they
are no longer the only point of comparison: Section~\ref{sec:res-models} places the tuned SVR against
\NModels{} models under equal conditions, and the finding that gradient-boosted trees lead on
this benchmark is connected to the broader literature in Section~\ref{sec:rw-benchmark}.

\pt{R2.3 Disclose LLM use.}
\reply Disclosed in full; see E12. I have also acted on the symptoms you identified rather
than only on the disclosure: the thin subsections have been consolidated, and every figure is
now regenerated from the result records together with its caption, which removes the
figure/caption mismatches.

\pt{R2.4 The training procedure is unclear --- 3-fold HPO, then 10-fold evaluation on what
data?}
\reply This was genuinely ambiguous in the earlier version, and I have separated the two
loops explicitly. Section~\ref{sec:ablation-protocol} now states that hyperparameter selection uses an inner
$k$-fold cross-validation \emph{on the training portion only}, and that performance is
estimated by an outer resampling scheme over held-out data that the inner loop never sees;
Section~\ref{sec:ablation-protocol} states which is used for each reported quantity. Where both selection and
hyperparameter search occur inside the outer folds, the estimate is labelled ``nested''
(Table~\ref{tab:validation}).

\pt{R2.5 Section 3 should follow Section 4 (previous version).}
\reply Done: the proposed workflow (now Section~\ref{sec:workflow}) precedes the experimental setup (now
Section~\ref{sec:setup}), so the proposal is stated before the data used to evaluate it.

\pt{R2.6 What is the regression equation mentioned on Page 7?}
\reply That sentence referred to nothing defined in the manuscript and has been removed. The
SVR decision function is now given once, in Section~\ref{sec:svr}, with every symbol defined.

\pt{R2.7 Figure titles duplicate captions; Figure~9's title and caption disagree; figures
referenced by title appear unreferenced.}
\reply All embedded figure titles have been removed, so the caption is the sole description.
Every figure is regenerated from its result record, and each is referenced in the text by
its \LaTeX\ label rather than by title, so the reference, the caption and the content are
guaranteed to correspond.

\pt{R2.8 Table captions use arabic numerals but the text uses Roman; format is not Springer
Nature style.}
\reply Corrected: all tables are numbered and referenced with arabic numerals in Springer
Nature style, using the \texttt{sn-jnl} class throughout.

\pt{R2.9 References use an incorrect format; academic papers should have DOIs.}
\reply The bibliography has been rebuilt in \texttt{sn-mathphys-num} style, and every one
of the 29 references now resolves to a persistent link. Twenty-one carry a DOI, including two
that were missing and have been added (Vapnik's monograph,
\texttt{10.1007/978-1-4757-2440-0}, and Shapley's 1953 chapter,
\texttt{10.1515/9781400881970-018}). The remaining eight have no DOI because their venues do
not register them --- five are \emph{Journal of Machine Learning Research} articles, two are
NeurIPS proceedings papers, and one is a technical report --- so each of those now carries the
publisher's stable URL instead, verified to resolve to the correct paper. No reference is now
without either a DOI or a permanent link.

\pt{R2.10 Section 4.1.2 is not self-contained --- define $n$, $\phi$, $b$, $\mathbf{w}$,
$\gamma$, $\alpha$, $y$, $\mathbf{x}$, $\|\mathbf{w}\|^2$.}
\reply Agreed --- a paper explicitly about SVR should not omit its formulation.
Section~\ref{sec:svr} now gives the primal objective, the $\varepsilon$-insensitive loss, the dual and
the kernel expansion, with every symbol you list defined at first use, and Table~\ref{tab:notation} collects
the full notation.

\pt{R2.11 Table 3 placement; Section 4.2.2 should precede 4.2.1; results in Section 4 belong
in Results.}
\reply Section~\ref{sec:setup} has been restructured as you describe. The derived features are now listed
before the table that summarises them; the ordering of the two subsections is reversed; and
every empirical value that appeared in the old Section~\ref{sec:setup} has been moved into Section~\ref{sec:results}, so
Section~\ref{sec:setup} states the setup and reports no results.

\pt{R2.12 Hyperparameter tuning: why discretise $C$ and $\varepsilon$ for a randomised
search? Why not a grid? Why tune only the RBF kernel? Why only 20 iterations if the search
takes 9 seconds? Table~\ref{tab:statistics} is unnecessary.}
\reply Every part of this was a fair criticism and all are now addressed empirically rather
than argued. $C$, $\varepsilon$ and $\gamma$ are now drawn from continuous log-uniform
distributions, so the randomised search is a genuine randomised search. Section~\ref{sec:res-kernels}
(Table~\ref{tab:search}) reports the exhaustive grid alongside randomised search at three budgets: the
\GridFits{}-fit grid reaches $\GridRtwo$ in $\GridSecs$~s, against $\RandTwentyRtwo$ at 20
iterations ($\RandTwentySecs$~s), $\RandSixtyRtwo$ at 60 ($\RandSixtySecs$~s) and
$\RandTwoHundredRtwo$ at 200 ($\RandTwoHundredSecs$~s). I should correct one thing I would
otherwise have written here: the search costs minutes, not the nine seconds quoted in the
previous version, and the exhaustive grid costs \GridSecsMin{}~minutes. That is still
affordable, so your point stands in full --- there was no computational case for the small
budget, the grid is now reported alongside the randomised search, and the limitation is
removed rather than defended. The returns are also not monotone in the budget: 60 iterations
give the best score of the four settings, and 200 iterations cost seven times as much to
arrive at slightly less. That is single-split sampling noise, and it is part of why the
significance testing is done over repeated splits. Every kernel is now tuned over its own hyperparameters, not only the RBF
(Table~\ref{tab:kernels}). The previous version's Table 5 has been deleted and its content folded into the text.

\pt{R2.13 Why is stratification needed in a regression problem?}
\reply It is not needed, and the manuscript now says so explicitly rather than implying
otherwise. Section~\ref{sec:eval-protocol} states that stratification is not required in
regression and that the deciles are used for one narrow purpose: to remove a source of
variation between configurations that are being compared against each other, so that a
difference between two cells reflects the configuration rather than a difference in the
target distribution of their partitions. Nothing in the paper's conclusions rests on it ---
the repeated-split analysis of Section~\ref{sec:res-stats}, which is where every significance
claim is made, draws fresh unstratified partitions, and the spatially blocked design of
Section~\ref{sec:res-validation} is the deliberate alternative where dependence between
observations genuinely matters. I have kept the device and justified it rather than removing
it, since removing it would add noise to the ablation comparisons without changing any
conclusion; I am happy to drop it entirely if you would prefer.

\pt{R2.14 Were the other algorithms also tuned? Did they use the same pipeline?}
\reply They are now, and they do. Section~\ref{sec:res-models} gives every model the same twelve features,
the same training rows, the same scaling stage and the same search budget of 20 randomised
iterations with three-fold cross-validation. Section~\ref{sec:fixed} states this explicitly. In the
earlier version the comparison models were used at their defaults, which --- as you imply ---
made the ranking claim uninterpretable.

\pt{R2.15 The ablation needs detail: did you start from Preethi et al.'s setting and add
components progressively? Was anything held fixed?}
\reply The design is now stated in full in Section~\ref{sec:ablation-protocol}, and it no longer proceeds by a single
progressive path. The baseline cell is the configuration described by Preethi et al.\ (no
scaling, raw features, defaults), at $R^2 = \AblBaseline$; from there all \AblNcells{} cells
of the four-factor design are evaluated, so no component's contribution depends on what was
held fixed while it was introduced. The best cell is \AblBestCell{} at $\AblBestRtwo$.

\pt{R2.16 Section 7.1 repeats earlier results; consider merging Sections 6 and 7.}
\reply The Discussion has been rewritten as three subsections that make claims rather than
restate numbers --- what the evidence supports, what it does not support, and threats to
validity --- with no repetition of the Results tables.

\pt{R2.17 If the search takes 9 seconds, the stated limitations do not hold --- why sample
3{,}000 instances, and why not increase the HPO budget? LIBSVM's complexity is not a reason
for a less ambitious setting.}
\reply You were right and I have not defended any of it. The 3{,}000-row subset is removed
from the main results (Table~\ref{tab:size} shows the full fit takes $\SizeFullSecs$~s), the search
budget is increased and justified empirically (Table~\ref{tab:search}), and the LIBSVM-complexity argument
has been deleted from the limitations. What remains in Section~\ref{sec:disc-threats} are limitations I can
actually defend: the spatial optimism, the census cap at \$500{,}001, the 1990 vintage of the
data, the within-domain nature of the additional datasets, and the reading error of roughly
$\pm 0.02$ in the values taken from Preethi et al.'s figures.

\subsection*{Minor issues}

\pt{R2.18 Add a reference for the hedonic pricing model (Section 1.1).}
\reply Added: Rosen (1974), with Limsombunchai (2004), Fan et al.\ (2006) and Park and Bae
(2015) for the machine-learning alternative.

\pt{R2.19 Add references for the California Housing dataset.}
\reply Added: Pace and Barry (1997) for the dataset's introduction and Pedregosa et al.\
(2011) for its scikit-learn distribution.

\pt{R2.20 Section 1.4's second sentence is confusing --- 10 new features but top 12 used ---
and contradicts the experiments.}
\reply The confusion was real: ten features are \emph{derived}, giving eighteen candidates
alongside the eight raw features, from which twelve are selected. The earlier phrasing
implied twelve were selected from ten. Section~\ref{sec:contributions} and Section~\ref{sec:derived} now state the counts
explicitly at each stage, and Figure~\ref{fig:pipeline} shows the dimensionality at each transition.

\pt{R2.21 Section 2.3's title should reflect its content.}
\reply Section~\ref{sec:related} has been restructured into three subsections with titles that state their
content: \emph{Feature scaling in kernel methods}, \emph{The California Housing benchmark},
and \emph{Comparing models and attributing gains}.

\pt{R2.22 Table 1 is unnecessary (previous version).}
\reply Deleted, with the small amount of information it carried folded into the text.

\pt{R2.23 Section 3.2 repeats ``characteristics'' in one sentence.}
\reply Rewritten.

\pt{R2.24 How can ``wide residual spread'' be observed in the caption of Figure~7?}
\reply It could not be, from that figure. The caption is rewritten to describe only what the
panels show --- residuals centred near zero with dispersion increasing at higher predicted
values, and the vertical band of true values at 5.0 produced by the census recording cap.

\section{Response to Reviewer 3}

Thank you for the recommendation of minor revision and for a report that identified the same
central problem as Reviewer~1 while also being precise about where the caveats belonged.

\pt{R3.1 The central comparison does not hold together: Stage~A yields $-0.054$ against
Preethi et al.'s $\approx 0.60$, and the gap is never explained.}
\reply This was the most important observation in the three reports, and pursuing it changed
the paper. The gap could not be explained because the premise was wrong: $0.60$ is the score
of their linear family, and their SVR is $\approx \PreethiReportedRtwo$. Rather than either
of the two options you offered, I have taken the third: the misreading is corrected and
stated openly (Sections~1.2 and~5.1), their configuration is reproduced directly, and the
reproduced values $\RepUnscaledDefault$ and $\RepUnscaledTuned$ bracket their actual result.
The ``explains the prior result'' language is now supported by a direct reproduction rather
than by an assumed correspondence.

\pt{R3.2 The training-set-size caveat must accompany the ranking claim wherever it appears,
not only in Limitations.}
\reply Rendered unnecessary rather than relocated: all \NModels{} models are now trained on
identical data (Section~\ref{sec:res-models}), so the ranking no longer carries an asymmetry that needs a
caveat. Both training sizes are reported, and the ranking is stated with its training size
attached in the abstract, in Section~\ref{sec:res-models} and in the conclusion.

\pt{R3.3 The Introduction asserts ``this does not restrict SVR as a learning algorithm''
before any evidence is presented.}
\reply Removed and replaced by the question being tested. Section~\ref{sec:questions} now poses four
research questions (Q1--Q4), each cross-referenced to the section that answers it, so the
Introduction states what is asked rather than what will be concluded.

\pt{R3.4 The claim that full-dataset training ``would be expected to yield lower RMSE'' is
untested.}
\reply It is now tested rather than asserted. Table~\ref{tab:size} fits the same tuned configuration at
three training sizes; RMSE falls from $\SizeThreeKrmse$ at 3{,}000 rows to $\SizeFullRmse$ on
the full training set, so the direction of the earlier claim is confirmed --- but it is now
a measurement.

\pt{R3.5 Soften ``confirms'' to ``supports'' or ``is consistent with''.}
\reply Done throughout. ``Confirms'' now appears only where a specific prediction was tested
directly and the outcome was as predicted, as in R3.4 above.

\vspace{10pt}
\hrule
\vspace{8pt}

\noindent I am grateful for the care evident in all three reports. The requirement to
reproduce the prior configuration directly is what uncovered my misreading of the reference
value, and the requirement to evaluate the full ablation design is what showed that my
headline attribution was an artefact of ordering. The paper is more modest than it was and, I
believe, considerably more solid.

\vspace{10pt}
\noindent Yours sincerely,\\[10pt]
Emmanuel Adutwum\\
Soka University of America\\
\texttt{eadutwum@soka.edu}

\end{document}
